\section{An explicit map with positive displacement}\label{app:auxiliary}

We give a direct construction of the map needed in
Lemma~\ref{lem:displacement}. This proves the Hilbert-space instance of
the Benyamini--Sternfeld theorem~\cite{BenyaminiSternfeld1983} directly.
The aim is to construct a Lipschitz map
$G:B\to H$ bounded away from zero and equal to the identity when
$\norm{x}\ge1/2$. Then $f(x)=-G(x)/\norm{G(x)}$ has the required
displacement. A tube about a curve with a bounded tail supplies $G$.

\begin{proposition}\label{aux:map}
Let $H=L^2([0,1];\R^2)$ be a real Hilbert space and let
$B=\{x\in H:\norm{x}\le1\}$. There is a globally Lipschitz map $f:B\to B$
such that
\[
  \norm{f(x)}=1,
  \qquad \norm{f(x)-x}\ge\frac12
  \quad(x\in B).
\]
\end{proposition}

Define functions on $\R$ by
\[
 a(t)=
 \begin{cases}
 t/8,&t\le2,\\
 5/16-(3-t)^2/16,&2<t<3,\\
 5/16,&t\ge3,
 \end{cases}
 \qquad
 \theta(t)=
 \begin{cases}
 0,&t\le1,\\
 (t-1)^2/2,&1<t<2,\\
 t-3/2,&t\ge2.
 \end{cases}
\]
For $w\in[0,1]$, put
\[
 v(t)(w)=(\cos(\theta(t)w),\sin(\theta(t)w)),
 \qquad \gamma(t)=a(t)v(t).
\]
Thus $\norm{v(t)}=1$, $\gamma(0)=0$, and $\gamma$ follows a straight line
for $t\le1$; its tail for $t\ge1$ lies in the ball of radius $5/16$.
The next lemma shows that, although the positive tail is bounded, curve
points whose parameters differ by a fixed amount remain uniformly
separated. This permits a tube of fixed radius along the whole curve.

\begin{lemma}\label{aux:tube}
The curve $\gamma$ is injective, has a bounded Lipschitz derivative, and
\[
 c:=\inf_{t\in\R}\norm{\gamma'(t)}=\frac18.
\]
Its unit tangent $u(t)=\gamma'(t)/\norm{\gamma'(t)}$ is Lipschitz,
$\ip{u(t)}{v(t)}\ge0$, and $u(t)=v(t)$ for $t\le1$.
There is a radius $0<\rho<1/32$ such that every point in
\[
 U=\{x\in H:\operatorname{dist}(x,\gamma(\R))<\rho\}
\]
has a unique nearest point $\gamma(t(x))$. Writing
$e(x)=x-\gamma(t(x))$, we have $e(x)\perp u(t(x))$, and
\begin{equation}\label{aux:parameter}
 |t(x)-t(y)|\le16\norm{x-y}
 \quad\text{if }x,y\in U\text{ and }\norm{x-y}<\rho.
\end{equation}
\end{lemma}

\begin{proof}
Let $J(\xi_1,\xi_2)=(-\xi_2,\xi_1)$. Since $0\le w\le1$, the
scalar trigonometric Taylor remainders are uniform in $w$, which
justifies differentiation in $H$. We obtain
\[
 \gamma'(t)(w)=a'(t)v(t)(w)+a(t)\theta'(t)wJv(t)(w),
 \qquad
 \norm{\gamma'(t)}^2=a'(t)^2+\frac{a(t)^2\theta'(t)^2}{3}.
\]
For $t\le2$ the first term is $1/64$; for $t\ge2$ we have
$a(t)\ge1/4$ and $\theta'(t)=1$. Equality in the claimed lower speed
bound holds for $t\le1$. The first derivatives match at the junctions
$1,2,3$, and their derivatives on the intervening intervals are bounded.
On the unbounded negative interval the angular derivatives vanish.
Consequently $\gamma'$ is bounded and Lipschitz. The same is true of
$v$ and $u$, and
$\ip{u(t)}{v(t)}=a'(t)/\norm{\gamma'(t)}\ge0$.

We need the following uniform separation property:
\begin{equation}\label{aux:separation}
 \sigma(d):=\inf_{|t-s|\ge d}\norm{\gamma(t)-\gamma(s)}>0
 \qquad(d>0).
\end{equation}
Indeed, with $\operatorname{sinc}(q)=\sin(q)/q$ for $q\ne0$ and
$\operatorname{sinc}(0)=1$, we have
\[
 \ip{v(s)}{v(t)}=\operatorname{sinc}(\theta(t)-\theta(s)).
\]
Since $|\operatorname{sinc}(q)|<1$ for $q\ne0$, distinct angular
parameters give noncollinear unit vectors. On $(-\infty,1]$ the
coefficient $a$ is strictly increasing, whereas $\theta$ is strictly
increasing on $(1,\infty)$ and $a$ is positive there. This proves
injectivity.

If \eqref{aux:separation} failed, choose $s_n<t_n$ with
$t_n-s_n\ge d$ and $\norm{\gamma(t_n)-\gamma(s_n)}\to0$, and pass to
subsequences for which both parameters have limits in the extended real
line. Finite limits contradict injectivity. If both limits are
$-\infty$, the distances equal $(t_n-s_n)/8$ eventually. If only
$s_n\to-\infty$, the first radii tend to infinity while the second
remain bounded. If $s_n\to s\in\R$ and $t_n\to+\infty$, the displayed
inner product tends to zero, and the squared distances tend to
$a(s)^2+(5/16)^2$. Finally, if both parameters tend to $+\infty$, then
eventually
\[
 \norm{\gamma(t_n)-\gamma(s_n)}^2
 =2(5/16)^2\bigl(1-\operatorname{sinc}(t_n-s_n)\bigr).
\]
This is bounded away from zero because
$\sup_{q\ge d}\operatorname{sinc}(q)<1$, by continuity and decay at
infinity. These cases establish \eqref{aux:separation}.

Write $K=\Lip(\gamma')$ and $K_u=\Lip(u)$. Choose $d>0$ with
$Kd\le c/2$. The estimate
\[
 \norm{\gamma(t)-\gamma(s)-(t-s)\gamma'(s)}
 \le\frac K2|t-s|^2
\]
then implies
\begin{equation}\label{aux:tangent}
 |\ip{\gamma(t)-\gamma(s)}{u(s)}|
 \ge\frac{3c}{4}|t-s|
 \qquad(|t-s|\le d).
\end{equation}
Choose $\rho>0$ so that
\begin{equation}\label{aux:radius}
 \rho<\frac1{32},\qquad 4\rho<\sigma(d),\qquad \rho K_u<\frac c4.
\end{equation}

For $x\in U$, take a minimizing sequence $\gamma(t_n)$ for the
distance from $x$ to the curve. Its late terms satisfy
$\norm{x-\gamma(t_n)}<2\rho$, so their pairwise distances are less
than $4\rho$. By \eqref{aux:separation} their parameters differ by
less than $d$. Fixing one late parameter therefore confines all later
ones to a compact real interval. A convergent subsequence attains the
minimum. Differentiating $\norm{x-\gamma(t)}^2$ at any minimizing
parameter gives $(x-\gamma(t))\perp u(t)$.

For $x,y\in U$ with $\norm{x-y}<\rho$, choose any minimizing
parameters $t,s$ and write $x=\gamma(t)+e_t$,
$y=\gamma(s)+e_s$. The curve points are less than $3\rho$ apart,
so $|t-s|<d$. Orthogonality and \eqref{aux:tangent} give
\begin{align*}
 \norm{x-y}
 &\ge |\ip{\gamma(t)-\gamma(s)}{u(s)}|
       -|\ip{e_t}{u(s)-u(t)}|\\
 &\ge\left(\frac{3c}{4}-\rho K_u\right)|t-s|
 \ge\frac c2|t-s|.
\end{align*}
Taking $y=x$ proves uniqueness; the same inequality proves
\eqref{aux:parameter}.
\end{proof}

We next arrange that a residual perpendicular to the curve's tangent
is carried to a vector perpendicular to $v(t)$. This will let us change
the coefficient along $v(t)$ without cancellation by that residual.

\begin{lemma}\label{aux:rotation}
There is a family of orthogonal operators $R_t:H\to H$, Lipschitz
in operator norm, such that $R_tu(t)=v(t)$ and $R_t$ is the identity
for $t\le1$.
\end{lemma}

\begin{proof}
Suppress $t$ and put $\kappa=\ip{u}{v}\ge0$. Define
\[
 Ah=v\ip{u}{h}-u\ip{v}{h},
 \qquad R=\mathrm{Id}+A+\frac{A^2}{1+\kappa}.
\]
If $u,v$ are independent, then $A$ is skew-adjoint,
$A^2=-(1-\kappa^2)\mathrm{Id}$ on their span, and $R$ restricts there
to $\kappa\mathrm{Id}+A$. This restriction is orthogonal and sends
$u$ to $v$, while $R$ is the identity on the orthogonal complement.
If $u=v$, then $A=0$ and $R=\mathrm{Id}$. These exhaust the cases
because $\kappa\ge0$. The formula, the Lipschitz dependence of
$u,v$, and $1+\kappa\ge1$ show that $t\mapsto R_t$ is Lipschitz
in operator norm, including where $u=v$.
\end{proof}

\begin{proof}[Proof of Proposition~\ref{aux:map}]
Fix the tube radius from Lemma~\ref{aux:tube} and the rotations from
Lemma~\ref{aux:rotation}. Set
\[
 b(t)=\min\{a(t),-1/8\},\qquad
 q(t)=b(t)-a(t)=-\max\{a(t)+1/8,0\}.
\]
Then $|q(t)|\le7/16$ and $\Lip(q)\le1/8$. Let
$\beta,\chi:[0,\infty)\to[0,1]$ be the continuous piecewise linear
functions specified by
\[
 \beta(r)=
 \begin{cases}
 1,&r\le\rho/4,\\
 2-4r/\rho,&\rho/4<r<\rho/2,\\
 0,&r\ge\rho/2,
 \end{cases}
 \qquad
 \chi(r)=
 \begin{cases}
 1,&r\le\rho/2,\\
 2-2r/\rho,&\rho/2<r<\rho,\\
 0,&r\ge\rho.
 \end{cases}
\]
Define $G:B\to H$ to equal the identity outside $U$. For $x\in B\cap U$,
write $t=t(x)$, $e=x-\gamma(t)$, $r=\norm{e}$, and set
\begin{equation}\label{aux:G}
 \begin{split}
 G(x)&=\gamma(t)+\beta(r)q(t)v(t)
       +\bigl((1-\chi(r))\mathrm{Id}+\chi(r)R_t\bigr)e\\
 &=x+\beta(r)q(t)v(t)+\chi(r)(R_t-\mathrm{Id})e.
 \end{split}
\end{equation}

First, $G$ is globally Lipschitz on $B$. For tube points at distance
less than $\rho$, \eqref{aux:parameter} controls their parameters,
and, with $M=\sup_t\norm{\gamma'(t)}$, their residuals satisfy
\[
 \norm{e(x)-e(y)}\le(1+16M)\norm{x-y}.
\]
Moreover, $r(x)=\operatorname{dist}(x,\gamma(\R))$ is
$1$-Lipschitz. Every factor in the second line of \eqref{aux:G} is
therefore Lipschitz for such pairs with uniform constants; the factors
are uniformly bounded, using $\norm{e}<\rho$, $|q|\le7/16$,
$\norm{v}=1$ and $\norm{R_t}=1$. The product estimates give one
finite Lipschitz bound for all these pairs. Also, on all of $B$,
\[
 \norm{G(x)}\le1+7/16+2\rho<3/2,
\]
so pairs at distance at least $\rho$ are controlled by $3/\rho$.
For the remaining pairs, let $x\in U$ and $y\notin U$. Both cutoffs
satisfy $\beta(r),\chi(r)\le2(\rho-r)/\rho$, whence
\[
 \norm{G(x)-x}
 \le\left(\frac{7}{8\rho}+4\right)(\rho-r).
\]
Since distance to the curve is $1$-Lipschitz and $y\notin U$,
$\rho-r\le\norm{x-y}$. Together with $G(y)=y$, this proves the
required bound across the tube boundary. Outside the tube $G$ is the
identity.

Next, we claim that
\begin{equation}\label{aux:nonzero}
 \norm{G(x)}\ge\rho/4\qquad(x\in B).
\end{equation}
Outside $U$, the fact that $0=\gamma(0)$ gives
$\norm{G(x)}=\norm{x}\ge\rho$. Inside $U$, if $r\le\rho/2$, then
$\chi(r)=1$ and $R_te\perp v(t)$, with $\norm{R_te}=r$. For
$r\le\rho/4$, the coefficient along $v(t)$ is $b(t)\le-1/8$;
for $\rho/4\le r\le\rho/2$, the perpendicular component has norm
at least $\rho/4$. If $\rho/2\le r<\rho$, then $\beta(r)=0$ and
the last term in the first line of \eqref{aux:G} has norm at most
$r<\rho$. When $t>1$, we have $a(t)\ge1/8$, so
$\norm{G(x)}\ge1/8-\rho>\rho/4$. When $t\le1$, we have
$R_t=\mathrm{Id}$ and $e\perp v(t)$, so
$\norm{G(x)}\ge r\ge\rho/2$. This proves \eqref{aux:nonzero}.

Finally, $G(x)=x$ whenever $\norm{x}\ge1/2$. This holds by definition
outside $U$. Inside $U$, the reverse triangle inequality gives
\[
 |a(t)|=\norm{\gamma(t)}\ge\norm{x}-r>1/2-\rho>5/16.
\]
Since $a(t)\le5/16$ everywhere, this forces $a(t)<-1/8$ and $t\le1$.
Thus $q(t)=0$ and $R_t=\mathrm{Id}$, so \eqref{aux:G} again gives
$G(x)=x$.

Define
\[
 f(x)=-\frac{G(x)}{\norm{G(x)}}.
\]
Normalization on vectors of norm at least $\rho/4$ is Lipschitz with
constant at most $8/\rho$, so $f$ is globally Lipschitz by
\eqref{aux:nonzero}, and its values have norm one. If
$\norm{x}<1/2$, then $\norm{f(x)-x}\ge1-\norm{x}>1/2$.
Otherwise $G(x)=x$, and
$\norm{f(x)-x}=1+\norm{x}\ge3/2$.
\end{proof}
